\clearpage\pagestyle{matSimple}\phantomsection\label{reperage}\pdfbookmark[0]{Repérage du programme}{bm-reperage}
\matPageTitre{Repérage du programme}{Grande Section · objectifs « à partir de 5 ans » du programme de maternelle}
{\fontsize{9.5}{13.5}\selectfont Chaque objectif d’apprentissage du programme pour cette tranche d’âge est relié aux activités qui le travaillent. Les numéros renvoient aux pages du cahier.\par}\vspace{5mm}
{\fontsize{9.5}{13}\selectfont\begin{tabularx}{\linewidth}{@{}Xr@{}}\toprule
{\otSemi Objectif du programme} & {\otSemi Activités}\\\midrule
Comprendre qu’une quantité ne dépend ni de la nature des objets ni de leur organisation (jusqu’à dix) & \hyperref[act-gs-1]{1}\\[1.6mm]
Un de plus, un de moins : le nombre suivant, le nombre précédent & \hyperref[act-gs-8]{8}, \hyperref[act-gs-9]{9}\\[1.6mm]
Parcourir une collection en passant une et une seule fois par chaque élément & \hyperref[act-gs-4]{4}\\[1.6mm]
Dénombrer une collection jusqu’à dix en utilisant des compositions & \hyperref[act-gs-3]{3}\\[1.6mm]
Constituer une collection d’un cardinal donné, y compris par son écriture chiffrée & \hyperref[act-gs-5]{5}\\[1.6mm]
Comparer des quantités et des nombres écrits en chiffres & \hyperref[act-gs-6]{6}, \hyperref[act-gs-7]{7}\\[1.6mm]
Composer et décomposer les nombres jusqu’à dix ; les doubles & \hyperref[act-gs-11]{11}, \hyperref[act-gs-12]{12}, \hyperref[act-gs-13]{13}, \hyperref[act-gs-14]{14}\\[1.6mm]
Surcompter à partir d’un nombre donné & \hyperref[act-gs-10]{10}\\[1.6mm]
Associer une quantité, le nom d’un nombre et son écriture chiffrée (jusqu’à dix) & \hyperref[act-gs-2]{2}\\[1.6mm]
Écrire en chiffres les nombres de un à dix & \hyperref[act-gs-15]{15}, \hyperref[act-gs-16]{16}\\[1.6mm]
Comptine numérique : jusqu’à trente, à rebours de dix, de deux en deux, à partir d’un nombre donné & \hyperref[act-gs-17]{17}, \hyperref[act-gs-18]{18}, \hyperref[act-gs-19]{19}, \hyperref[act-gs-20]{20}\\[1.6mm]
Rang : premier, dernier, deuxième, avant-dernier ; changer de point de départ & \hyperref[act-gs-21]{21}, \hyperref[act-gs-22]{22}, \hyperref[act-gs-23]{23}\\[1.6mm]
Déterminer l’effet d’un déplacement : avancer, reculer & \hyperref[act-gs-26]{26}\\[1.6mm]
Construire la bande numérique jusqu’à dix & \hyperref[act-gs-24]{24}, \hyperref[act-gs-25]{25}\\[1.6mm]
Problèmes de parties-tout, d’ajout et de retrait & \hyperref[act-gs-27]{27}, \hyperref[act-gs-28]{28}, \hyperref[act-gs-29]{29}\\[1.6mm]
Problèmes d’écart : « de plus », « de moins » & \hyperref[act-gs-30]{30}, \hyperref[act-gs-31]{31}\\[1.6mm]
Problèmes de position : retrouver la case de départ & \hyperref[act-gs-32]{32}\\[1.6mm]
Problèmes de groupements et de partage (avec un reste) & \hyperref[act-gs-33]{33}, \hyperref[act-gs-34]{34}, \hyperref[act-gs-35]{35}\\[1.6mm]
Solides : décrire le cube et la pyramide par leurs faces & \hyperref[act-gs-38]{38}, \hyperref[act-gs-39]{39}\\[1.6mm]
Formes planes : carré, rectangle, triangle, disque, dans toutes les orientations & \hyperref[act-gs-36]{36}, \hyperref[act-gs-37]{37}\\[1.6mm]
Reproduire des assemblages de formes (au plus huit) & \hyperref[act-gs-40]{40}\\[1.6mm]
Utiliser la règle pour tracer & \hyperref[act-gs-41]{41}\\[1.6mm]
Longueurs : bande témoin, ranger jusqu’à cinq objets, reproduire une longueur & \hyperref[act-gs-42]{42}, \hyperref[act-gs-43]{43}, \hyperref[act-gs-44]{44}\\[1.6mm]
Masses : ordonner trois objets, réaliser l’équilibre & \hyperref[act-gs-45]{45}, \hyperref[act-gs-46]{46}\\[1.6mm]
Motifs : motifs évolutifs, transcription, structure commune, création & \hyperref[act-gs-47]{47}, \hyperref[act-gs-48]{48}, \hyperref[act-gs-49]{49}, \hyperref[act-gs-50]{50}, \hyperref[act-gs-51]{51}\\[1.6mm]
Se repérer : gauche et droite, codage de déplacements, quadrillage & \hyperref[act-gs-52]{52}, \hyperref[act-gs-53]{53}, \hyperref[act-gs-54]{54}, \hyperref[act-gs-55]{55}, \hyperref[act-gs-56]{56}\\[1.6mm]
\bottomrule\end{tabularx}\par}
